\subsection{Links to get started}
\label{links}

Overview to easily get started (high-level motivation and initial introductions on
youtube)\footnote{URLs partly edited for improved readability. Don't use the printed URL
directly, but rather click on the link to go to the linked target.}:
\begin{itemize}
    \scriptsize
    \item
    \href{https://www.youtube.com/watch?v=1cRFfYQYGxE&list=PLcGKfGEEONaBNsY_bOj8IhbCPvj6OAdWo&index=33}{\textit{https://www.thestrangeloop.com}}
    (Jack Rusher, Funny intro to GA with historical perspective)
    \item
    \href{https://www.youtube.com/watch?v=60z_hpEAtD8&t=17s}{\textit{https://www.youtube.com/@sudgylacmoe}}
    (sudgylacmoe, A swift intro to GA)
    \item
    \href{https://www.youtube.com/watch?v=qJHFTMF_pPk&list=PLffJUy1BnWj3deu0cqpk5CiePiwORDiCe}{\textit{https://www.youtube.com/@AllAngles/GA}}
    (All Angles - GA Intro)
    \item
    \href{https://www.youtube.com/watch?v=zgi-13F2Kec&list=PLsSPBzvBkYjxV8QMlOTYgrqa-nmlt_L7A&index=11}{\textit{https://www.youtube.com/@bivector/GAME23}}
    (De Keninck and Roelfs, GAME 2023 - GA Intro)
\end{itemize}

Sources for GA and PGA as a mathematical topic:
\begin{itemize}
    \scriptsize
    \item
    \href{https://projectivegeometricalgebra.org}{\textit{https://projectivegeometricalgebra.org}}
    (Eric Lengyel)
    \item
    \href{https://terathon.com/blog/poor-foundations-ga.html}{\textit{https://terathon.com/blog/poor-foundations-ga.html}}
    (Eric Lengyel)
    \item
    \href{https://bivector.net}{\textit{https://bivector.net}}
    (Leo Dorst et.al.)
    \item
    \href{https://www.youtube.com/@bivector}{\textit{https://www.youtube.com/@bivector}}
    (Youtube channel bivector)
    \item
    \href{https://en.wikipedia.org/wiki/Bivector}{\textit{https://en.wikipedia.org/wiki/Bivector}}
    (Wikipedia Bivector)
    \item
    \href{https://www.youtube.com/@sudgylacmoe}{\textit{https://www.youtube.com/@sudgylacmoe}}
    (Youtube channel sudgylacmoe)
    \item
    \href{https://www.youtube.com/watch?v=ItGlUbFBFfc}{\textit{https://www.youtube.com/@hestenes}}
    (Hestenes, Tutorial Geometric Calculus (relevant formulas))
    \item
    \href{https://en.wikipedia.org/wiki/Plane-based_geometric_algebra}{\textit{https://en.wikipedia.org/PGA}}
    (Plane-based GA in $R^*(3,0,1)$ instead of $R(3,0,1)$ as used here)
\end{itemize}

Sources for GA and PGA in more specific applications:
\begin{itemize}
    \scriptsize
    \item
    \href{https://www.youtube.com/watch?v=v-WG02ILMXA&t=1476s}{\textit{https://www.youtube.com/DynamicsPGA}}
    (Leo Dorst, Dynamics in PGA)
    \item
    \href{https://en.wikipedia.org/wiki/Newton–Euler_equations}{\textit{https://en.wikipedia.org/wiki/Newton-Euler}}
    (Newton-Euler equations)
    \item
    \href{https://en.wikipedia.org/wiki/Euler%27s_equations_(rigid_body_dynamics)}{\textit{https://en.wikipedia.org/wiki/Rigidbody}}
    (Rigid body dynamics)
    \item
    \href{https://en.wikipedia.org/wiki/Moment_of_inertia}{\textit{https://en.wikipedia.org/wiki/Inertia}}
    (Moment of inertia)
    \item
    \href{https://en.wikipedia.org/wiki/Screw_theory#Wrench}{\textit{https://en.wikipedia.org/wiki/ScrewTheory}}
    (Screw theory, twist and wrench)
    \item
    \href{https://en.wikipedia.org/wiki/Chasles%27_theorem_(kinematics)}{{https://en.wikipedia.org/wiki/Chasles}}
    (Chasles' theorem)
    \item
    \href{http://web.cs.iastate.edu/~cs577/handouts/plucker-coordinates.pdf}{{http://web.cs.iastate.edu/Pluecker}}
    (Plücker coordinates in projective spaces)

\end{itemize}

\newpage

% literature and references
\subsection{Literature}
\label{literature}

%\addcontentsline{toc}{section}{\numberline{}Literature}

\bibliographystyle{apalike}
\bibliography{9_ga_literature-db}